\chapter{Biological and Statistical Background}

\section{Biological Framework}

Cancer growth can be thought of as a dynamic process that unfolds within a population of cells. A tumor usually starts with a mutation of one cell, which subsequently starts dividing. During the subsequent cell division events, the initial mutant cell produces a lineage of cells, which grow into a subclone of descendants. During the growth of this lineage, cells accumulate somatic mutations and these are passed onto descendants. Similar to other evolutionary phenomena, somatic mutations bring diversity to the system, and the outcome depends on how various clones are fitter than others, expanding into large populations, staying rare, or dying off.

A tumor is an inherently stochastic process. Genetic mutations occur spontaneously either during the process of DNA replication or following any DNA damages and imprecise repair mechanism. Moreover, there is no guarantee that a potentially successful mutation will make a cell live longer or multiply: a cell carrying the mutation might simply die or remain small. Thus, when observing the final composition of a tumor, it is impossible to disentangle what exactly happened to individual cells in terms of biology and stochastic factors linked to cell proliferation and death, and mutation acquisition and extinction.

In this setting, mutations are often classified as passengers or drivers. Passenger mutations are neutral mutations that do not change cell biology and do not affect its proliferation or fitness. Yet, they are very useful as ancestry indicators: cells with identical passenger mutations share ancestry. Driver mutations represent changes that affect cell biology, such as growth, replication, survival, or other functions relevant to cancer evolution. These mutations may provide some selective advantage to a mutated cell, driving the formation of a subclone.

As a tumor grows, its evolutionary history can be represented with a phylogenetic tree. The tree is rooted at the first tumor cell that emerged and contains information about accumulated mutations and the formation of clones. Here, clones correspond to groups of cells that share identical sets of mutations, and a subclone appears in the case when a cell within a clone acquires another set of mutations and grows. In particular, at the moment of sequencing, only those cells that survive are observed, and it is required to infer the evolutionary process behind a given set of observed genotypes.

Another fundamental parameter in cancer is mutation rate. This rate can be defined, for instance, as the frequency of mutations appearing per cell division. Under some stochastic assumptions a specific number of mutations can be used as a proxy to measure branch length: the more mutations there are along a branch of a phylogenetic tree, the older the event, and vice versa. However, mutations do not necessarily accumulate uniformly along a branch.

Another type of somatic events in tumors is copy-number alterations. While a normal diploid region consists of a copy of two alleles, one from the mother and another from the father, cancer genomes can have gains, losses, or duplications of certain genomic fragments. In this work, we consider three main types of copy numbers, illustrated in Figure~\ref{fig:cna-states}: (2:1), (2:2), and (2:0). Copy-number (2:1) denotes trisomy, which means that one allele is in double copy while the other allele is only in one copy, making the total number of copies equal to three. Copy-number (2:2) denotes tetrasomy, which could result from WGD where both alleles are duplicated, making the total number of copies four. Finally, copy-number (2:0) denotes copy-neutral loss of heterozygosity (CNLOH), where an allele gets deleted while the remaining allele gets duplicated, resulting in total copies of two but one allele being completely deleted.

Another fundamental concept is clone growth rate. Growth rate defines how fast a specific clone grows over time. A faster-growing clone is able to reach a higher frequency of its occurrence within the tumor population. As a result, the final frequencies of clones depend not only on their age but also on their growth rates. Specifically, early clones have had more time to grow, yet their growth rates could be relatively slow. On the other hand, newer clones are younger, yet they can have higher growth rates and therefore be present in higher share. 

It becomes clear now that the analysis of tumor evolution requires probabilistic models to account for multiple aspects of the process. Number of mutations gives information about the timing of different events, clone proportions in the sampled cells give information about their expansion dynamics, and phylogenetic trees show ordering between the clones. While the constructed phylogenetic tree should not be treated as exact ground truth, it remains an effective way to model the tumor process.
% --- colori e helper ---
\definecolor{premut}{HTML}{E0A92B}
\definecolor{postmut}{HTML}{2FA36B}
\colorlet{alleleA}{blue!60}
\colorlet{alleleB}{red!65}
\newcommand{\allelebar}[3]{\draw[fill=#3,draw=#3!65!black,line width=.4pt,rounded corners=2pt] (#1,#2) rectangle ++(0.5,2.0);}
\newcommand{\muttri}[3]{\fill[#3] (#1-0.18,#2+0.12) -- (#1-0.18,#2-0.12) -- (#1,#2) -- cycle;}
\newcommand{\timetri}[2]{\fill[#2] (#1-0.12,2.66) -- (#1+0.12,2.66) -- (#1,2.44) -- cycle;}

\begin{figure}[H]\centering
\resizebox{\textwidth}{!}{%
\begin{tikzpicture}[x=1cm,y=1cm]
% bounding box simmetrico attorno a x=8.5 (centro dei blocchi) -> figura centrata
\useasboundingbox (-0.6,0.3) rectangle (17.6,11.4);
\node[font=\bfseries] at (8.38,11.0) {Copy-number states \& mutations};

% ===== ORIGIN 1:1 =====
\node[font=\small,text=gray] at (8.38,10.3) {Diploid origin - \textbf{1:1}};
\allelebar{7.75}{8.0}{alleleA}
\allelebar{8.75}{8.0}{alleleB}
\muttri{7.75}{9.65}{premut}
\muttri{8.75}{9.25}{premut}

\draw[-{Latex[length=2mm]},gray] (8.1,7.95) -- (3.2,6.9);
\draw[-{Latex[length=2mm]},gray] (8.5,7.95) -- (8.5,6.9);
\draw[-{Latex[length=2mm]},gray] (8.9,7.95) -- (13.8,6.9);

% ===== 2:1  (A,A,B) =====
\allelebar{1.9}{4.8}{alleleA}\allelebar{2.95}{4.8}{alleleA}\allelebar{4.0}{4.8}{alleleB}
\muttri{1.9}{6.45}{premut}\muttri{2.95}{6.45}{premut}\muttri{4.0}{6.05}{premut}
\muttri{1.9}{5.35}{postmut}\muttri{2.95}{5.20}{postmut}\muttri{4.0}{5.45}{postmut}
\node[font=\bfseries] at (3.2,4.5){2:1};
\node[font=\small,text=gray] at (3.2,3.9){Trisomy};

% ===== 2:2  (A,A,B,B) =====
\allelebar{6.675}{4.8}{alleleA}\allelebar{7.725}{4.8}{alleleA}\allelebar{8.775}{4.8}{alleleB}\allelebar{9.825}{4.8}{alleleB}
\muttri{6.675}{6.45}{premut}\muttri{7.725}{6.45}{premut}\muttri{8.775}{6.05}{premut}\muttri{9.825}{6.05}{premut}
\muttri{6.675}{5.35}{postmut}\muttri{7.725}{5.20}{postmut}\muttri{8.775}{5.45}{postmut}\muttri{9.825}{5.25}{postmut}
\node[font=\bfseries] at (8.55,4.5){2:2};
\node[font=\small,text=gray] at (8.25,3.9){Whole genome doubling};


% ===== 2:0  (A,A, B perso) =====
\allelebar{12.5}{4.8}{alleleA}\allelebar{13.55}{4.8}{alleleA}
\draw[draw=gray,line width=.4pt,rounded corners=2pt,dashed,fill=gray!8] (14.6,4.8) rectangle ++(0.5,2.0);
\draw[gray,line width=.4pt] (14.6,4.8)--++(0.5,2.0);\draw[gray,line width=.4pt] (15.1,4.8)--(14.6,6.8);
\muttri{12.5}{6.45}{premut}\muttri{13.55}{6.45}{premut}
\fill[premut,opacity=0.3] (14.6-0.18,6.05+0.12)--(14.6-0.18,6.05-0.12)--(14.6,6.05)--cycle;
\muttri{12.5}{5.35}{postmut}\muttri{13.55}{5.20}{postmut}
\node[font=\bfseries] at (13.8,4.5){2:0};
\node[font=\small,text=gray] at (13.8,3.9){Copy neutral loss of heterozygosity};


% ===== PSEUDO-TIME =====
\def\xTimeStart{1.66}
\def\xTimeEnd{15.5}
\def\xTimeMid{8.6}
\def\yTime{2.40}

\renewcommand{\timetri}[2]{%
  \fill[#2] (#1-0.11,\yTime+0.22) -- (#1+0.11,\yTime+0.22) -- (#1,\yTime+0.04) -- cycle;
}

\draw[line width=.5pt] (\xTimeStart,\yTime) -- (\xTimeEnd,\yTime);

\draw[line width=.5pt] (\xTimeStart,\yTime-0.08) -- (\xTimeStart,\yTime+0.08);
\draw[line width=.5pt] (\xTimeEnd,\yTime-0.08) -- (\xTimeEnd,\yTime+0.08);

\node[font=\small,text=gray] at (\xTimeStart-0.15,\yTime-0.30) {0};
\node[font=\small,text=gray] at (\xTimeEnd+0.15,\yTime-0.30) {1};

\node[font=\small,text=gray] at (\xTimeMid,\yTime-0.58) {Pseudo-time};

\def\xTCN{7.80}
\draw[dashed,gray] (\xTCN,\yTime-0.08) -- (\xTCN,\yTime+0.58);
\node[font=\footnotesize,text=gray] at (\xTCN,\yTime+0.74) {$t_{\mathrm{CN}}$};

\timetri{3.74}{premut}
\timetri{5.70}{premut}
\timetri{9.90}{postmut}
\timetri{11.58}{postmut}
\timetri{13.40}{postmut}

% ===== LEGEND (con positioning) =====
\begin{scope}
% primo simbolo
\node[
  draw=alleleA!65!black,
  fill=alleleA,
  rounded corners=1pt,
  minimum width=0.32cm,
  minimum height=0.32cm,
  inner sep=0pt
] (legA) at (0.95,0.91) {};

\node[
  anchor=west,
  right=2mm of legA,
  font=\small,
  text=gray
] (txtA) {Allele A};

% secondo simbolo
\node[
  draw=alleleB!65!black,
  fill=alleleB,
  rounded corners=1pt,
  minimum width=0.32cm,
  minimum height=0.32cm,
  inner sep=0pt,
  right=10mm of txtA
] (legB) {};

\node[
  anchor=west,
  right=2mm of legB,
  font=\small,
  text=gray
] (txtB) {Allele B};

% triangolo premut
\node[
  regular polygon,
  regular polygon sides=3,
  shape border rotate=90,
  draw=premut,
  fill=premut,
  minimum size=0.32cm,
  inner sep=0pt,
  right=10mm of txtB
] (legPre) {};

\node[
  anchor=west,
  right=2mm of legPre,
  font=\small,
  text=gray
] (txtPre) {Pre-CNA ($\alpha$)};

% triangolo postmut
\node[
  regular polygon,
  regular polygon sides=3,
  shape border rotate=90,
  draw=postmut,
  fill=postmut,
  minimum size=0.32cm,
  inner sep=0pt,
  right=10mm of txtPre
] (legPost) {};

\node[
  anchor=west,
  right=2mm of legPost,
  font=\small,
  text=gray
] (txtPost) {Post-CNA ($\beta$)};

% lost copy
\node[
  draw=gray,
  dashed,
  fill=gray!8,
  rounded corners=1pt,
  minimum width=0.32cm,
  minimum height=0.32cm,
  inner sep=0pt,
  right=10mm of txtPost
] (legLost) {};

\node[
  anchor=west,
  right=2mm of legLost,
  font=\small,
  text=gray
] {Lost copy};
\end{scope}
\end{tikzpicture}%
}
\captionsetup{font=footnotesize}
\caption{Top: copy-number states $2{:}1$, $2{:}2$, $2{:}0$ from a diploid $1{:}1$ origin. Pre-CNA mutations (gold, $\alpha$) sit at a fixed locus on every surviving copy of their allele; post-CNA mutations (green, $\beta$) are independent per copy. In $2{:}0$ the lost allele (dashed) removes its mutation. Bottom: pseudo-time, with $\alpha$ before and $\beta$ after the CNA time $t_{\mathrm{CN}}$.}
\label{fig:cna-states}
\end{figure}
\subsection{Sequencing}

Sequencing measures the genetic makeup of a tumour and provides information about
somatic mutations, copy-number changes, and clonality. Two sequencing strategies are
relevant to this work: bulk sequencing and single-cell sequencing.

Bulk sequencing extracts DNA from the whole cell population, so the signal it returns
is averaged over all the sampled cells. For each genomic site one observes the
\emph{variant allele frequency} (VAF), the fraction of sequencing reads at that site
that carry the mutation \parencite{cocomello2025timing}. Bulk sequencing is cheap and
widely used, but the averaging hides which cells carry which mutations, so
reconstructing the clonal composition from VAFs alone is difficult.

Single-cell sequencing measures cells one at a time. This makes it possible to see
which mutations co-occur in the same cell, to assign cells to clones, and to
reconstruct the phylogenetic tree. The gain in resolution is what makes single-cell
data so valuable for studying tumour evolution. The price is noise: because each cell
contains very little DNA, it must be amplified, and amplification introduces allelic
dropout, false positives, and uneven, sometimes insufficient, sequencing depth. As a
result, not finding a mutation in a given cell does not mean the mutation is absent.

Single-cell sequencing is the basis of this work: the assignment of cells to clones
and the number of mutations along each branch of the tree give us the information we
need about event timings and growth rates. The limitations described above concern
\emph{accuracy}, the noise in the measured genotypes, and not the tree topology
itself. A separate issue is \emph{sampling}: because only part of the tumour is
sequenced, the reconstructed phylogeny, and in particular which clone appears as the
most recent common ancestor (MRCA), depends on which cells are sampled and need not
match the true evolutionary history. Figure~\ref{fig:sampling-mrca} illustrates how
different samples of the same tumour can yield different reconstructed trees and a
different MRCA.


% --- colori e helper (invariati) ---
\definecolor{clA}{HTML}{8A8F98}\definecolor{clB}{HTML}{4C78C8}\definecolor{clC}{HTML}{F58518}
\definecolor{clD}{HTML}{54A24B}\definecolor{clE}{HTML}{B279A2}\definecolor{clF}{HTML}{E45756}
\definecolor{mrcacol}{HTML}{E0A92B}
\newcommand{\cells}[2]{\fill[black!75] (#1-0.18,#2-0.6) circle(.6mm);\fill[black!75] (#1,#2-0.64) circle(.6mm);\fill[black!75] (#1+0.18,#2-0.6) circle(.6mm);}
\newcommand{\cn}[4]{\filldraw[fill=#3,draw=#3!55!black,line width=.4pt] (#1,#2) circle(3.5mm);\node[text=white,font=\bfseries\small] at (#1,#2){#4};}
\newcommand{\sn}[4]{\filldraw[fill=#3,draw=black,line width=1.2pt] (#1,#2) circle(3.5mm);\node[text=white,font=\bfseries\small] at (#1,#2){#4};\cells{#1}{#2}}
\newcommand{\fn}[4]{\filldraw[fill=#3,fill opacity=0.22,draw=gray,dashed,line width=.4pt] (#1,#2) circle(3.5mm);\node[text=white,opacity=0.45,font=\bfseries\small] at (#1,#2){#4};}
\newcommand{\eg}[4]{\draw[line width=.7pt,black!50] (#1,#2)--(#3,#4);}
\newcommand{\fe}[4]{\draw[line width=.7pt,black!50,opacity=0.22,dashed] (#1,#2)--(#3,#4);}
\newcommand{\mr}[2]{\draw[mrcacol,line width=1.2pt] (#1,#2) circle(5.5mm);}

\begin{figure}[H]\centering
\resizebox{\textwidth}{!}{%
\begin{tikzpicture}[x=1cm,y=1cm]
\node[font=\bfseries] at (6.0,13.5){How sampling changes the reconstructed MRCA};


% ===================== TL : TRUE TREE =====================
\eg{2.0}{11.5}{0.9}{10.3}\eg{2.0}{11.5}{2.9}{10.3}\eg{2.9}{10.3}{1.8}{9.1}\eg{2.9}{10.3}{3.7}{9.1}\eg{3.7}{9.1}{2.9}{7.9}
\cn{2.0}{11.5}{clA}{A}\cn{0.9}{10.3}{clC}{C}\cn{2.9}{10.3}{clB}{B}\cn{1.8}{9.1}{clE}{E}\cn{3.7}{9.1}{clD}{D}\cn{2.9}{7.9}{clF}{F}
\node[font=\small,text=gray] at (2.3,12.5){True tree (6 clones)};

% ===================== TR : SAMPLING 1 {C,F} -> A =====================
\node[font=\small,text=gray] at (9.3,12.5){Sampling 1: $\{C,F\}$ - MRCA = A};
\draw[mrcacol,dashed,line width=.6pt] (7.3,11.5)--(11.2,11.5);
\node[font=\small,text=mrcacol!75!black,anchor=east] at (7.4,11.5){MRCA = A};
\eg{9.0}{11.5}{7.9}{10.3}\eg{9.0}{11.5}{9.9}{10.3}\eg{9.9}{10.3}{10.7}{9.1}\eg{10.7}{9.1}{9.9}{7.9}
\fe{9.9}{10.3}{8.8}{9.1}
\cn{9.0}{11.5}{clA}{A}\sn{7.9}{10.3}{clC}{C}\cn{9.9}{10.3}{clB}{B}\fn{8.8}{9.1}{clE}{E}\cn{10.7}{9.1}{clD}{D}\sn{9.9}{7.9}{clF}{F}
\mr{9.0}{11.5}

% ===================== BL : SAMPLING 2 {E,F} -> B =====================
\begin{scope}[yshift=-0.8cm]
\draw[mrcacol,dashed,line width=.6pt] (0.4,4.6)--(4.3,4.6);
\node[font=\small,text=mrcacol!75!black,anchor=west] at (4.5,4.6){MRCA = B};
\eg{2.9}{4.6}{1.8}{3.4}\eg{2.9}{4.6}{3.7}{3.4}\eg{3.7}{3.4}{2.9}{2.2}
\fe{2.0}{5.8}{0.9}{4.6}\fe{2.0}{5.8}{2.9}{4.6}
\fn{2.0}{5.8}{clA}{A}\fn{0.9}{4.6}{clC}{C}\cn{2.9}{4.6}{clB}{B}\sn{1.8}{3.4}{clE}{E}\cn{3.7}{3.4}{clD}{D}\sn{2.9}{2.2}{clF}{F}
\mr{2.9}{4.6}
\node[font=\small,text=gray] at (2.3,6.7){Sampling 2: $\{E,F\}$ - MRCA = B};

% ===================== BR : SAMPLING 3 {D,F} -> D =====================
\node[font=\small,text=gray] at (9.3, 6.7){Sampling 3: $\{D,F\}$ - MRCA = D};
\draw[mrcacol,dashed,line width=.6pt] (7.4,3.4)--(11.3,3.4);
\node[font=\small,text=mrcacol!75!black,anchor=east] at (7.4,3.4){MRCA = D};
\eg{10.7}{3.4}{9.9}{2.2}
\fe{9.0}{5.8}{7.9}{4.6}\fe{9.0}{5.8}{9.9}{4.6}\fe{9.9}{4.6}{8.8}{3.4}\fe{9.9}{4.6}{10.7}{3.4}
\fn{9.0}{5.8}{clA}{A}\fn{7.9}{4.6}{clC}{C}\fn{9.9}{4.6}{clB}{B}\fn{8.8}{3.4}{clE}{E}\cn{10.7}{3.4}{clD}{D}\sn{9.9}{2.2}{clF}{F}
\mr{10.7}{3.4}


% ===================== LEGEND (distanziata) =====================

\node[circle,fill=clA,draw=black,line width=1.1pt,minimum size=4.8mm,inner sep=0pt] (lgA) at (0.5,0.05){};
\node[font=\small,text=gray,right=2mm of lgA] (lgAt){Sampled clone};
\node[circle,fill=clA,fill opacity=0.22,draw=gray,dashed,line width=.4pt,minimum size=4.8mm,inner sep=0pt,right=9mm of lgAt] (lgB){};
\node[font=\small,text=gray,right=2mm of lgB] (lgBt){Not reconstructed};
\node[circle,draw=mrcacol,line width=1.1pt,minimum size=4.8mm,inner sep=0pt,right=9mm of lgBt] (lgC){};
\node[font=\small,text=gray,right=2mm of lgC]{MRCA};
\end{scope}
\end{tikzpicture}
}
\captionsetup{font=footnotesize}
\caption{Examples of how different samplings lead to different reconstructed trees}
\label{fig:sampling-mrca}
\end{figure}

\subsection{Real Data and ProCESS}

Validation of the proposed tumour evolution model on authentic single-cell data is clearly the ultimate aim because the latter encapsulates the complexity of the biological process. However, several issues make such validation problematic. Datasets of single-cell sequencing are expensive to obtain, hard to access, and many publicly available cancer data were not collected with the aim of reconstructing clonal evolution, often lacking the mutation numbers, clone assignments or a phylogenetic tree needed for our model. More importantly, there is no ground truth: even after inference, sequencing-based phylogeny cannot be considered to be accurate, and timings of the evolutionary events and the growth rates of clones are always unknown. Hence, the validation must rely on simulation.

The first step of the validation will be conducted assuming that the sequence was generated according to the process defined by the model. Inability to reproduce the generating parameter values would mean that the inference algorithm either fails in some particular cases or there is no enough information contained in the dataset.

Another more stringent validation procedure will be carried out based on the Programmable Cancer Evolution Spatial Simulator (ProCESS). As opposed to simplified simulation methods used for data generation, ProCESS produces realistic cancer histories while keeping track of all process parameters and allowing us to compare them directly with our inferred estimates. Moreover, ProCESS enables single-cell sampling from the tumour in question, allowing us to create a partially and possibly biased sample similar to sequencing data. Additionally, simulations under the same parameter settings can be repeated multiple times for evaluating stability of the inference results.

Hence, ProCESS provides a compromise between the idealistic simulation and real data as it maintains biologically relevant details, yet retains perfect control over the data-generating process and ground truth information.

\subsection{The Inference Problem}

The goal of this thesis can be stated in one sentence: from a single snapshot of
single-cell sequencing data, reconstruct \emph{when} the events in a tumour's history
happened and \emph{how fast} each clone grew. For a tumour made of $K$ clones, this
means estimating the birth time of every clone, the time of every copy-number
alteration, and the growth rate of every clone, each reported with its uncertainty
rather than as a single point estimate.

We start from a specific set of inputs. The first is the reconstructed phylogenetic
tree over the clones, obtained from the single-cell data by an aforementioned method. What
the tree provides are the \emph{parental relations}: for each clone it tells us the
parent clone from which it descends. These relations are the backbone of the problem,
because they fix the \emph{order} of the events even when their exact times are
unknown: a clone is always younger than its parent, and a copy-number alteration on a
branch must have occurred after the parent was born and before the child split off.
Writing $t_i$ for the birth time of clone $i$ on a scale where the tumour starts at
$0$ and sequencing happens at $1$, the tree thus imposes
%
\begin{equation}
    0 < t_{\mathrm{pa}(i)} < t_i < 1 .
\end{equation}
%
We do not treat the tree as exact ground truth; we take its structure, who descends
from whom, as the scaffold on which the timing is inferred.

On top of this scaffold, three kinds of measurement carry the information we exploit:
the number of somatic mutations along each branch, which acts as a clock for how long
the branch lasted; for the regions affected by a copy-number alteration, the split of
those mutations into the ones that arose before and after the event, which locates the
alteration within its branch; and the number of cells assigned to each clone, which
reflects how much the clone expanded before sequencing.

These signals do not pin down the timings on their own. The number of mutations and
the size of a clone both depend on the combination of how fast it grew and how long it
lived, so a clone that grew quickly for a short time can look much like one that grew
slowly for longer. Separating the timing of an event from the growth of the clone is
therefore the central difficulty of the problem. The model developed in Chapter~2
addresses it by combining these sources of evidence in a single Bayesian framework and
by sharing growth information between related clones along the tree, so that timings
and growth rates are recovered together, with calibrated uncertainty.
\section{Statistical Framework}

\subsection{Brief introduction to bayesian inference}
Bayesian inference is a consistent mathematical formalism for learning about the 
underlying quantities from observations. While the frequentist framework treats 
parameters as fixed but unknown constants, the Bayesian one regards parameters as 
random variables and makes use of Bayes' rule to update a prior distribution in 
light of observed data.

Bayes' theorem, the foundation of Bayesian inference, relates the posterior 
probability distribution $p(\theta \mid \mathbf{y})$, representing our beliefs 
about unknown parameters $\theta$ given data $\mathbf{y}$, the likelihood 
$p(\mathbf{y} \mid \theta)$, reflecting our belief about how likely the observed 
data were generated under different $\theta$ values, and the prior $p(\theta)$:
%
\begin{equation}
    p(\theta \mid \mathbf{y}) = \frac{p(\mathbf{y} \mid \theta)\, p(\theta)}{p(\mathbf{y})},
\end{equation}
%
where $p(\mathbf{y}) = \int p(\mathbf{y} \mid \theta)\, p(\theta)\, \mathrm{d}\theta$ 
is called the marginal likelihood (also known as evidence or partition function) and 
serves as a normalising factor. Since the marginal likelihood does not depend on the 
value of the unknown parameters, the posterior distribution may equivalently be 
written as
%
\begin{equation}
    p(\theta \mid \mathbf{y}) \propto p(\mathbf{y} \mid \theta)\, p(\theta).
\end{equation}
%
The above equation captures the essence of Bayesian inference, namely that the 
posterior is proportional to the product of the likelihood and the prior.

An important feature of the Bayesian framework is that it gives rise to different 
forms of priors. If there is sufficient prior knowledge about the parameters, one 
can use an informative prior and concentrate prior probability mass around plausible 
parameter values. Weakly informative priors act as very mild regularisers with 
little effect on the posterior distribution. On the other hand, in many problems, 
data alone is insufficient to recover the parameter values; that is, there is no 
unique solution to the problem. In such problems, the prior acts as an additional 
source of information, ensuring that the resulting posterior is well-regularised. 
The latter scenario is encountered in the current thesis because mutation counts
and clone abundances only constrain the product $w_i \, \Delta t_i$, not its two
factors separately: a clone that grew quickly for a short time looks the same as
one that grew slowly for longer. It is precisely because the data cannot separate
$w_i$ from $\Delta t_i$ that we introduce a prior, which supplies the information
the likelihood is missing along this direction and yields a well-defined posterior.

One generalisation of the above framework is the hierarchical model. In this setup, 
the parameters themselves are treated as random variables whose distribution depends 
on some additional hyperparameters. More formally, let $\theta_1, \ldots, \theta_K$ 
be parameters whose values are distributed according to
%
\begin{equation}
    \theta_i \mid \phi \sim p(\theta_i \mid \phi), \quad i \in \{1, \ldots, K\}.
\end{equation}
%
Here, $\phi$ denotes hyperparameters, and we impose a hyperprior $p(\phi)$ on $\phi$. 
It follows that the joint posterior distribution is given by
%
\begin{equation}
    p(\theta_1, \ldots, \theta_K, \phi \mid \mathbf{y}) \propto 
    p(\mathbf{y} \mid \theta_1, \ldots, \theta_K)\, 
    \prod_{i=1}^{K} p(\theta_i \mid \phi)\, p(\phi).
\end{equation}
%
One of the key advantages of the hierarchical approach is partial pooling, where 
information from one unit is pooled into the rest of the units, making it possible 
for poorly informed units to borrow strength from better-informed ones. In the model 
proposed in this thesis, a hierarchical prior is imposed on clone-specific growth 
rates based on a random walk on the log scale along phylogenetic branches.

In practice, the posterior distribution $p(\theta \mid \mathbf{y})$ is rarely 
obtainable analytically, so we need to rely on some approximation method. The 
leading methodology used in applied Bayesian inference is Markov chain Monte Carlo 
(MCMC), which builds a Markov chain whose limiting distribution coincides with the 
desired posterior. Posterior quantities are estimated from post-burn-in Markov chain 
samples. An efficient MCMC scheme is Hamiltonian Monte Carlo (HMC) and its adaptive 
version, the No-U-Turn Sampler (NUTS). Both use gradient information from the 
log-posterior to generate distant jumps, ensuring that the chain mixes quickly even 
in high-dimensional or highly correlated spaces. The models proposed in this thesis 
are coded in Stan, a probabilistic programming language that automatically computes 
gradients and implements NUTS as its default sampler.

\subsection{Hamiltonian Monte Carlo}
Hamiltonian Monte Carlo (HMC) represents one approach to Markov chain Monte Carlo 
sampling inspired by classical mechanics. This involves defining a new momentum 
variable $\rho$, which must have the same dimensions as $\theta$, according to
%
\begin{equation}
    H(\theta, \rho) = -\log p(\theta \mid \mathbf{y}) + \frac{1}{2}\, \rho^{\top} M^{-1} \rho,
\end{equation}
%
with $M$ representing the mass matrix, normally set to the identity matrix. The 
first term can be viewed as potential energy and the second as kinetic energy. 
Sampling proceeds through Hamiltonian dynamics generated by numerically integrating 
the equations of motion
%
\begin{align}
    \frac{\mathrm{d}\theta}{\mathrm{d}s} &= M^{-1} \rho, \\[4pt]
    \frac{\mathrm{d}\rho}{\mathrm{d}s} &= \nabla_{\!\theta} \log p(\theta \mid \mathbf{y}),
\end{align}
%
using a leapfrog integrator with integration step size $\varepsilon$ and $L$ 
leapfrog steps. A Metropolis correction is applied to the result to account for 
any errors in the integration. In essence, the leapfrog algorithm makes use of 
gradients of the log-posterior to generate distant proposals, which are rarely 
rejected, resulting in chains exploring the posterior distribution better than 
random walk samplers.

One of the major issues with Hamiltonian Monte Carlo is setting the hyperparameters 
$\varepsilon$ and $L$ appropriately. The No-U-Turn Sampler (NUTS) eliminates this 
problem by automatically computing the value of $L$, increasing it until a reversal 
in the leapfrog trajectory occurs, explaining its name. Step size tuning is 
performed automatically as part of a burn-in period. NUTS is the default sampler 
used by Stan, and the gradients $\nabla_{\!\theta} \log p(\theta \mid \mathbf{y})$ 
are automatically calculated using automatic differentiation.